% ID: SEL-D0051-OB-E01
% ESTADO: validada. S052.
\textbf{a)} Sobre $y=0$ queda $2\le x\le4$. En $x=2$, la cota superior es $y=(48-2)/8=23/4$. Al resolver simultáneamente $5x-4y=20$ y $x+8y=48$ se obtiene $(8,5)$. Los cuatro vértices cumplen todas las restricciones.
\begin{center}
\setlength{\unitlength}{0.9mm}
\begin{picture}(91,51)(-6,-8)
\qbezier(0,0)(41,0)(82,0)
\qbezier(0,0)(0,20)(0,40)
\put(83,0){$x$}
\put(1,40){$y$}
\qbezier(16,0)(24,0)(32,0)
\qbezier(32,0)(48,15)(64,30)
\qbezier(64,30)(40,32.25)(16,34.5)
\qbezier(16,34.5)(16,17.25)(16,0)
\put(11,-6){\scriptsize $(2,0)$}
\put(28,-6){\scriptsize $(4,0)$}
\put(65,30){\scriptsize $(8,5)$}
\put(12,36){\scriptsize $(2,23/4)$}
\end{picture}
\end{center}
\textbf{b)} Los valores de $F$ en $(2,0),(4,0),(8,5),(2,23/4)$ son, respectivamente, $4,8,76,73$. Por linealidad, el mínimo es $\boxed{4}$ en $(2,0)$ y el máximo es $\boxed{76}$ en $(8,5)$. Cada extremo es único porque sólo uno de los vértices alcanza su valor.

\textbf{c)} Como $100>76$, ningún punto del recinto cumple $F(x,y)=100$.
